\section{Graph and warehouse model}
\label{sec:graph}

\begin{lede}
The warehouse is a grid of open and blocked cells. This section defines how far
apart two cells are, which cells are corridors, which are junctions, and which
single cells the whole map depends on.
\end{lede}

\subsection{Grid graph}
\label{sec:graph-grid}

\code{GridGraph} is a 4-connected grid: from $v = (r, c)$, the neighbour set is
\begin{equation}
  N(v) = \{(r \pm 1, c),\, (r, c \pm 1)\} \cap V,
\end{equation}
with the four deltas fixed in that order so that iteration is deterministic.
Degree is $\deg(v) = |N(v)|$, and $\Delta(G) = \max_v \deg(v)$
(\code{GridGraph.max\_degree}) appears in the PIBT complexity bound
(\Cref{sec:pibt-complexity}).

\subsection{BFS distance maps}
\label{sec:graph-distances}

For a goal $g$, \code{compute\_bfs\_distance\_map(g)} runs a reverse BFS from
$g$ and returns, for every vertex $v$,
\begin{equation}
  \Dg{g}(v) = d_{\mathrm{route}}(v, g)
  \qquad \text{(hop count; } \INF \text{ if unreachable)},
\end{equation}
in $O(|V|)$ since the grid is unweighted. The map is cached per goal and is
the single most reused primitive in the codebase: it backs route planning,
progress scoring, congestion lookahead, task-assignment cost, and the
A$^\ast$ heuristic for both the direction-aware router (\Cref{sec:routing}) and
the baseline planners (\Cref{sec:baselines}) --- all without recomputation.
\code{Warehouse.precompute()} eagerly builds $\Dg{g}$ for every pickup,
delivery and parking vertex.

\code{shortest\_route(source, goal, horizon)} walks $\Dg{g}$ greedily: at each
step move to the neighbour $n$ minimising $\Dg{g}(n)$, ties broken by the fixed
neighbour order, so the result is deterministic. An optional \code{horizon}
truncates the route after that many steps, used when only a lookahead prefix is
needed (\Cref{sec:congestion-downstream}).

\subsection{Structural features}
\label{sec:graph-structure}

Computed once by an iterative Hopcroft--Tarjan lowlink DFS
(\code{\_compute\_articulation\_and\_bridges}, non-recursive to avoid Python's
recursion limit on large maps):

\begin{itemize}
\item \textbf{Articulation points}: vertices whose removal disconnects $G$.
\item \textbf{Bridges}: edges whose removal disconnects $G$ --- equivalently,
  edges lying on no cycle.
\item \textbf{Dead ends}: $\{v : \deg(v) \leq 1\}$.
\item \textbf{Intersections}: $\{v : \deg(v) \geq 3\}$ --- exactly the vertices
  that separate one aisle from the next (\Cref{sec:graph-aisles}).
\end{itemize}

\code{satisfies\_pibt\_reachability()} returns
$(\text{dead ends} = \emptyset) \wedge (\text{bridges} = \emptyset)$, a rough
check of PIBT's sufficient condition that every edge lie on a simple cycle ---
the condition its progress guarantee assumes, and the one
\Cref{sec:contribution-soft} shows a hard one-way rule violates.

\code{bfs\_ball(center, radius)} returns every vertex within \code{radius} hops
of a centre, used for deadlock clustering (\Cref{sec:deadlock-detection}) and
escape-vertex search (\Cref{sec:recovery-levels}).

\subsection{Aisle segmentation}
\label{sec:graph-aisles}

An \textbf{aisle} is a maximal straight run of the graph after removing all
intersection vertices --- the corridors \emph{between} junctions. Each aisle's
\code{vertices} list is ordered from one endpoint to the other
(\code{\_order\_component}: find a local-degree-$\leq 1$ endpoint, then walk
neighbour to neighbour). Moving to a higher index is \textsc{forward}, to a
lower index \textsc{reverse}, and
\begin{equation}
  \mathrm{entry\_direction}(v) =
  \begin{cases}
    \textsc{forward} & \text{if } \mathrm{index}(v) \leq (\ell - 1)/2,\\
    \textsc{reverse} & \text{otherwise,}
  \end{cases}
\end{equation}
for an aisle of length $\ell$: entering nearer the start means heading forward
through it.

\begin{caveat}
\textbf{Segmentation cuts at every turn.} Segmenting by connected component
alone --- the literal reading --- produces L-, U- and T-shaped ``aisles'':
\variant{warehouse\_corridors} yielded two 25-cell U shapes each spanning a
whole corridor plus both vertical links. \textsc{forward} has no single compass
meaning on such a shape, and a one-way rule blocks its corner in both senses.
\code{Aisle.axis} reports \code{row} or \code{col}, and is empty only for a
single-cell aisle.
\end{caveat}

\paragraph{Capacity.} \code{\_aisle\_capacity($\ell$)} has two models
(\param{aisle\_capacity\_model}):
\begin{align}
  \text{\code{"length"}:}\quad
    &\mathrm{cap} = \mathrm{clamp}\!\left(\lceil \rho \ell \rceil,\ 1,\ \param{aisle\_capacity}\right),\\
  \text{\code{"throughput"}:}\quad
    &\mathrm{cap} = \mathrm{clamp}\!\left(\lceil \rho \ell / T_{\min} \rceil,\ 1,\ \param{aisle\_capacity}\right),
\end{align}
with $\rho = \param{aisle\_capacity\_ratio}$ (default $1.0$) and
$T_{\min} = \param{minimum\_aisle\_lock\_time}$ (default 20). The rationale for
the second: a single-file corridor is throttled by how fast the robot at its
\emph{exit} can leave, not by how many robots physically fit inside. Packing it
to $\ell$ just builds a queue that cannot drain.

Each aisle's minimum lock is
$\max(2, \min(T_{\min}, 2\ell))$, capped so a short aisle does not lock
disproportionately long for its own size.

\paragraph{Manageability.} An aisle is permanently \textsc{open} --- never
given a direction --- if either $\ell < \param{directional\_aisle\_min\_length}$
(default 4: too short for a rule to be worth anything) or some edge of the
aisle, or either boundary edge to an adjacent intersection, is a bridge.
One-waying a cut edge disconnects the graph, which no amount of pricing
repairs.
